% GENERATED FILE. Edit canonical lessons, then run npm run generate:books.
% canonical-source-hash: fnv1a64:4e6ba7578e3014c6
% canonical-lessons: GU-C394-hamal, GU-C394-pataro, GU-C394-potlu, GU-C394-sarkavvu, GU-C394-jhalvu

\chapter{Porters and Trunks}
\label{ch:gu-porters-and-trunks}
\hlchaptermodality{394}

\begin{chapteropening}
\textbf{By the end of this chapter:} \emph{I can ask a porter to carry my luggage.}

Complete the last lesson of chapter 394: I can ask a porter to carry my luggage.
\end{chapteropening}

\section[\texorpdfstring{\gu{હમાલ} (hamāl)}{હમાલ (hamāl)}]{\gu{હમાલ} (hamāl) — a porter}

\label{lesson:GU-C394-hamal}



\begin{quote}
Before the new one: say the Gujarati for a whistle, then the Gujarati for to couple.
\end{quote}

\subsection*{You'll want to know: \gu{હમાલ}}
\textbf{\gu{હમાલ}} — \emph{hamāl} — \textquotedblleft{}a porter\textquotedblright{}.

\begin{quote}
\gu{હમાલે} \gu{મારો} \gu{સામાન} \gu{ટ્રેનમાં} \gu{મૂક્યો}.
\end{quote}

\begin{quote}
\emph{hamāle māro sāmān ṭrenmā̃ mūkyo.}
\end{quote}

\begin{quote}
The porter put my luggage on the train.
\end{quote}

A masculine noun.

\begin{cousinweb}
\textbf{\gu{હમાલ}} (\emph{hamāl}) is usually traced to Arabic \emph{ḥammāl}, \textquotedblleft{}a carrier of loads\textquotedblright{}, which reached Gujarati through Persian; at a big station the porters carry bags on their heads.
\end{cousinweb}

\begin{grammarlens}[title={What you've built}]
One more word for getting around.
\end{grammarlens}

\subsection*{Your turn}
\begin{itemize}
\raggedright
  \item \emph{Say it:} \emph{hamāl}
  \item \emph{Say it:} \emph{hamāl}, once more
  \item \emph{Say it:} \emph{hamāle māro sāmān ṭrenmā̃ mūkyo.} — The porter put my luggage on the train.
  \item \emph{Recall:} say the Gujarati for a whistle, then the Gujarati for to couple, then say \emph{hamāl} again
\end{itemize}

\begin{tcolorbox}[breakable,colback=teal!4,colframe=teal!35!black,title={Before you move on}]
What does \gu{હમાલ} mean? (\textquotedblleft{}A porter\textquotedblright{}.) Say it once more.
\end{tcolorbox}
\section[\texorpdfstring{\gu{પટારો} (paṭāro)}{પટારો (paṭāro)}]{\gu{પટારો} (paṭāro) — a trunk (a large box)}

\label{lesson:GU-C394-pataro}



\begin{quote}
Before the new one: say the Gujarati for to couple, then the Gujarati for a porter.
\end{quote}

\subsection*{You'll want to know: \gu{પટારો}}
\textbf{\gu{પટારો}} — \emph{paṭāro} — \textquotedblleft{}a trunk (a large box)\textquotedblright{}.

\begin{quote}
\gu{દાદીનો} \gu{જૂનો} \gu{પટારો} \gu{ઘણો} \gu{ભારે} \gu{છે}.
\end{quote}

\begin{quote}
\emph{dādīno jūno paṭāro ghaṇo bhāre che.}
\end{quote}

\begin{quote}
Grandmother's old trunk is very heavy.
\end{quote}

A masculine noun.

\begin{cousinweb}
\textbf{\gu{પટારો}} (\emph{paṭāro}) goes back to Sanskrit \emph{piṭaka}, \textquotedblleft{}a basket, a box\textquotedblright{}; CDIAL lists the Gujarati word for a large box (CDIAL 8164).
\end{cousinweb}

\begin{grammarlens}[title={What you've built}]
One more word for getting around.
\end{grammarlens}

\subsection*{Your turn}
\begin{itemize}
\raggedright
  \item \emph{Say it:} \emph{paṭāro}
  \item \emph{Say it:} \emph{paṭāro}, once more
  \item \emph{Say it:} \emph{dādīno jūno paṭāro ghaṇo bhāre che.} — Grandmother's old trunk is very heavy.
  \item \emph{Recall:} say the Gujarati for to couple, then the Gujarati for a porter, then say \emph{paṭāro} again
\end{itemize}

\begin{tcolorbox}[breakable,colback=teal!4,colframe=teal!35!black,title={Before you move on}]
What does \gu{પટારો} mean? (\textquotedblleft{}A trunk (a large box)\textquotedblright{}.) Say it once more.
\end{tcolorbox}
\section[\texorpdfstring{\gu{પોટલું} (poṭlũ)}{પોટલું (poṭlũ)}]{\gu{પોટલું} (poṭlũ) — a bundle (tied in cloth)}

\label{lesson:GU-C394-potlu}



\begin{quote}
Before the new one: say the Gujarati for a porter, then the Gujarati for a trunk (a large box).
\end{quote}

\subsection*{You'll want to know: \gu{પોટલું}}
\textbf{\gu{પોટલું}} — \emph{poṭlũ} — \textquotedblleft{}a bundle (tied in cloth)\textquotedblright{}.

\begin{quote}
\gu{તેણે} \gu{કપડાંનું} \gu{પોટલું} \gu{માથા} \gu{પર} \gu{મૂક્યું}.
\end{quote}

\begin{quote}
\emph{teṇe kapḍā̃nũ poṭlũ māthā par mūkyũ.}
\end{quote}

\begin{quote}
He put the bundle of clothes on his head.
\end{quote}

A neuter noun.

\begin{cousinweb}
\textbf{\gu{પોટલું}} (\emph{poṭlũ}) goes back to \emph{*pōṭṭa}, \textquotedblleft{}a bundle\textquotedblright{}, a form CDIAL reconstructs beside Sanskrit \emph{pōṭṭala}; CDIAL lists the Gujarati word (CDIAL 8396).
\end{cousinweb}

\begin{grammarlens}[title={What you've built}]
One more word for getting around.
\end{grammarlens}

\subsection*{Your turn}
\begin{itemize}
\raggedright
  \item \emph{Say it:} \emph{poṭlũ}
  \item \emph{Say it:} \emph{poṭlũ}, once more
  \item \emph{Say it:} \emph{teṇe kapḍā̃nũ poṭlũ māthā par mūkyũ.} — He put the bundle of clothes on his head.
  \item \emph{Recall:} say the Gujarati for a porter, then the Gujarati for a trunk (a large box), then say \emph{poṭlũ} again
\end{itemize}

\begin{tcolorbox}[breakable,colback=teal!4,colframe=teal!35!black,title={Before you move on}]
What does \gu{પોટલું} mean? (\textquotedblleft{}A bundle (tied in cloth)\textquotedblright{}.) Say it once more.
\end{tcolorbox}
\section[\texorpdfstring{\gu{સરકાવવું} (sarkāvvũ)}{સરકાવવું (sarkāvvũ)}]{\gu{સરકાવવું} (sarkāvvũ) — to slide (something); to push along}

\label{lesson:GU-C394-sarkavvu}



\begin{quote}
Before the new one: say the Gujarati for a trunk (a large box), then the Gujarati for a bundle (tied in cloth).
\end{quote}

\subsection*{You'll want to know: \gu{સરકાવવું}}
\textbf{\gu{સરકાવવું}} — \emph{sarkāvvũ} — \textquotedblleft{}to slide (something); to push along\textquotedblright{}.

\begin{quote}
\gu{મેં} \gu{ભારે} \gu{થેલી} \gu{બારણા} \gu{પાસે} \gu{સરકાવી}.
\end{quote}

\begin{quote}
\emph{mẽ bhāre thelī bārṇā pāse sarkāvī.}
\end{quote}

\begin{quote}
I slid the heavy bag over to the door.
\end{quote}

Past, with \gu{તેણે} (he or she …): \gu{તેણે} \gu{પટારો} \gu{સરકાવ્યો}, \emph{teṇe paṭāro sarkāvyo}.

\begin{cousinweb}
\textbf{\gu{સરકાવવું}} (\emph{sarkāvvũ}) is the causative partner of \gu{સરકવું}, \textquotedblleft{}to slide\textquotedblright{}, which CDIAL lists under Sanskrit \emph{sárati}, \textquotedblleft{}runs, flows, glides\textquotedblright{} (CDIAL 13250): to make a thing slide.
\end{cousinweb}

\begin{grammarlens}[title={What you've built}]
One more word for getting around.
\end{grammarlens}

\subsection*{Your turn}
\begin{itemize}
\raggedright
  \item \emph{Say it:} \emph{sarkāvvũ}
  \item \emph{Say it:} \emph{sarkāvvũ}, once more
  \item \emph{Say it:} \emph{mẽ bhāre thelī bārṇā pāse sarkāvī.} — I slid the heavy bag over to the door.
  \item \emph{Recall:} say the Gujarati for a trunk (a large box), then the Gujarati for a bundle (tied in cloth), then say \emph{sarkāvvũ} again
\end{itemize}

\begin{tcolorbox}[breakable,colback=teal!4,colframe=teal!35!black,title={Before you move on}]
What does \gu{સરકાવવું} mean? (\textquotedblleft{}To slide (something); to push along\textquotedblright{}.) Say it once more.
\end{tcolorbox}
\section[\texorpdfstring{\gu{ઝાલવું} (jhālvũ)}{ઝાલવું (jhālvũ)}]{\gu{ઝાલવું} (jhālvũ) — to catch hold of; to grip}

\label{lesson:GU-C394-jhalvu}



\begin{quote}
Before the new one: say the Gujarati for a bundle (tied in cloth), then the Gujarati for to slide (something).
\end{quote}

\subsection*{You'll want to know: \gu{ઝાલવું}}
\textbf{\gu{ઝાલવું}} — \emph{jhālvũ} — \textquotedblleft{}to catch hold of; to grip\textquotedblright{}.

\begin{quote}
\gu{ભીડમાં} \gu{મીરાએ} \gu{બાળકનો} \gu{હાથ} \gu{ઝાલ્યો}.
\end{quote}

\begin{quote}
\emph{bhīḍmā̃ Mīrāe bāḷakno hāth jhālyo.}
\end{quote}

\begin{quote}
In the crowd Mira caught hold of the child's hand.
\end{quote}

Past, with \gu{તેણે} (he or she …): \gu{તેણે} \gu{હાથ} \gu{ઝાલ્યો}, \emph{teṇe hāth jhālyo}.

\begin{cousinweb}
\textbf{\gu{ઝાલવું}} (\emph{jhālvũ}) goes back to \emph{*jhal}, \textquotedblleft{}a sudden movement\textquotedblright{}, a form CDIAL reconstructs; CDIAL lists the Gujarati verb, \textquotedblleft{}to catch hold of\textquotedblright{} (CDIAL 5351): a quick grab.
\end{cousinweb}

\begin{grammarlens}[title={What you've built}]
One more word for getting around.
\end{grammarlens}

\subsection*{Your turn}
\begin{itemize}
\raggedright
  \item \emph{Say it:} \emph{jhālvũ}
  \item \emph{Say it:} \emph{jhālvũ}, once more
  \item \emph{Say it:} \emph{bhīḍmā̃ Mīrāe bāḷakno hāth jhālyo.} — In the crowd Mira caught hold of the child's hand.
  \item \emph{Recall:} say the Gujarati for a bundle (tied in cloth), then the Gujarati for to slide (something), then say \emph{jhālvũ} again
\end{itemize}

\begin{tcolorbox}[breakable,colback=teal!4,colframe=teal!35!black,title={Before you move on}]
What does \gu{ઝાલવું} mean? (\textquotedblleft{}To catch hold of; to grip\textquotedblright{}.) Say it once more.
\end{tcolorbox}
